%========================================
% [5] Buddhābhigīti
%========================================
% title-en: Chant in Praise of the Buddha
\section*{e5 | Buddhābhigīti}

\begin{chantsubsection}
  \chantnote
    {\{Handa mayaṃ buddhābhigītiṃ karomase\}}
  \chantnotetrans
    {\{Now let us chant in praise of the Buddha.\}}
\end{chantsubsection}

% ---- Opening verse ----
\begin{chantsubsection}
  \chantline
    {Buddhavārahantavaratādiguṇābhiyutto}
    {Endowed with the noble qualities of the Buddha, the Arahant, and the Excellent Ones,}
  \chantline
    {Suddhābhiñāṇakaruṇāhi samāgatatto}
    {his being complete with pure insight and compassion,}
  \chantline
    {Bodhesi yo sujanataṃ kamalaṃ va sūro}
    {he awakened the good people, like the sun [awakens] the lotus,}
  \chantline
    {Vandāmahaṃ taṃ araṇaṃ sirasā jinendaṃ}
    {I bow with my head to that peaceful one, the Lord of Conquerors.}
\end{chantsubsection}

% ---- Main praise (two pādas per line) ----
\begin{chantsubsection}
  \chantlinepair
    {Buddho yo sabbapāṇīnaṃ}
    {Saraṇaṃ khemamuttamaṃ}
    {The Buddha is the secure, supreme refuge for all living beings.}

  \chantlinepair
    {Paṭhamānussatiṭṭhānaṃ}
    {Vandāmi taṃ sirenāhaṃ}
    {As the first object of recollection, I bow to him with my head.}

  % footnote 1: use \footnotemark inside minipage text, then \footnotetext right after
  \chantlinepair
    {Buddhassāhaṃ dāso\footnotemark{} va}
    {Buddho me sāmikissaro}
    {I am the Buddha's servant; the Buddha is my lord and master.}
  \footnotetext{For ladies: change \textit{dāso} to \textit{dāsī}.}

  \chantlinepair
    {Buddho dukkhassa ghātā ca}
    {Vidhātā ca hitassa me}
    {The Buddha destroys suffering and brings about my welfare.}

  \chantlinepair
    {Buddhassāhaṃ niyyādemi}
    {Sarīrañjīvitañcidaṃ}
    {To the Buddha I dedicate this body and this very life.}

  % footnote 2
  \chantlinepair
    {Vandantohaṃ\footnotemark{} carissāmi}
    {Buddhasseva subodhitaṃ}
    {Paying homage, I will live (practice) in accordance with the Buddha’s clearly realized Dhamma.}
  \footnotetext{For ladies: change \textit{Vandantohaṃ} to \textit{Vandantīhaṃ}.}

  \chantlinepair
    {Natthi me saraṇaṃ aññaṃ}
    {Buddho me saraṇaṃ varaṃ}
    {I have no other refuge; the Buddha is my supreme refuge.}

  \chantlinepair
    {Etena saccavajjena}
    {Vaḍḍheyyaṃ satthu sāsane}
    {By this truthful utterance, may I grow in the Teacher’s dispensation.}

  % footnote 3
  \chantlinepair
    {Buddhaṃ me vandamānena\footnotemark{}}
    {Yaṃ puññaṃ pasutaṃ idha}
    {By my paying homage to the Buddha, whatever merit is produced here,}
  \footnotetext{For ladies: change \textit{vandamānena} to \textit{vandamānāya}.}

  \chantlinepair
    {Sabbepi antarāyā me}
    {Māhesuṃ tassa tejasā}
    {may all obstacles not afflict me through the power of his radiance.}
\end{chantsubsection}

\newpage

\begin{chantsubsection}
  \chantnote{\textit{--- bow, chanting softly ---}}
\end{chantsubsection}

% ---- Confession verse ----
\begin{chantsubsection}
  \chantline
    {Kāyena vācāya va cetasā vā}
    {By body, speech, or mind,}
  \chantline
    {Buddhe kukammaṃ pakataṃ mayā yaṃ}
    {whatever wrong I have done toward the Buddha,}
  \chantline
    {Buddho paṭiggaṇhatu accayantaṃ}
    {may the Buddha accept my confession,}
  \chantline
    {Kālantare saṃvarituṃ va buddhe}
    {so that in the future I may restrain myself with regard to the Buddha.}
\end{chantsubsection}
